% Generated draft: requires mathematical and source review.
\documentclass[11pt]{article}
\usepackage[margin=25mm]{geometry}
\usepackage{fontspec}
\setmainfont{Latin Modern Roman}
\usepackage{amsmath,amssymb,mathtools,unicode-math}
\setmathfont{Latin Modern Math}
\usepackage{enumitem,needspace,xurl,hyperref}
\setlength{\parindent}{0pt}
\setlength{\parskip}{5pt}
\setlength{\emergencystretch}{3em}
\newcommand{\N}{\mathbb{N}}
\newcommand{\Z}{\mathbb{Z}}
\newcommand{\Q}{\mathbb{Q}}
\newcommand{\R}{\mathbb{R}}
\newcommand{\C}{\mathbb{C}}
\newcommand{\F}{\mathbb{F}}
\newcommand{\A}{\mathbb{A}}
\newcommand{\PP}{\mathbb{P}}
\newcommand{\E}{\mathbb{E}}
\newcommand{\eps}{\varepsilon}
\newcommand{\dd}{\,\mathrm{d}}
\begin{document}
\section*{UnsolvedMath ID: 11000216}
\textbf{Source problem code:} catalogue entry
\subsection{Definitions and mathematical statement}
Retain the strata \(\mathcal H(\alpha)\) and \(\mathcal Q(\beta)\). In a usual convention, \(\mathcal H(\alpha)\) parametrizes holomorphic abelian differentials with zero orders \(\alpha\), and \(\mathcal Q(\beta)\) parametrizes quadratic differentials with prescribed orders \(\beta\), with at most simple poles in the finite-area setting. The excerpt does not specify genus, labels, connected components, or pole conventions. The group \(\operatorname{GL}^+(2,\R)\) acts by applying its matrices to the flat-coordinate charts.

\emph{Problem 9 (orbit closures for moduli spaces).} Determine, for every point in these strata, the closure of its \(\operatorname{GL}^+(2,\R)\)-orbit. Are the closures always complex-analytic orbifolds? Are they complex-algebraic? Characterize them geometrically. The closure topology and precise embedded/immersed orbifold convention must be fixed in the source setup. The analytic and algebraic structural questions and the complete geometric-classification request are distinct parts.

The source attributes to Kontsevich that a \(\operatorname{GL}^+(2,\R)\)-invariant complex-analytic subvariety is affine in period coordinates. Its global versus local meaning is not given; local affine equations do not themselves list all possible orbit closures.

The source also defines the area-one loci \(\mathcal H_1(\alpha)\subset\mathcal H(\alpha)\) and \(\mathcal Q_1(\beta)\). The subscript denotes area, not genus. It describes \(\mathcal H_1(\alpha)\) as a real-codimension-one, \(\operatorname{SL}_2(\R)\)-invariant level set defined in period coordinates by the quadratic Riemann bilinear relation, called the ``unit hyperboloid'' and locally modeled on a paraboloid. It says the stratum measure induces a natural invariant measure there and credits Masur and Veech with finite total measure for both area-one loci. The measure normalization and which action preserves the measure, rather than merely its class, are not specified; full \(\operatorname{GL}^+\) scaling must be distinguished from the area-preserving subgroup.

The registry marks the compound item solved. Its summary credits Eskin--Mirzakhani--Mohammadi with affine invariant orbit-closure structure, and Wright, Filip, and others with further progress, while explicitly retaining complete classification of affine invariant submanifolds as an active/open program. These supplied structural-resolution claims do not mean that every geometric closure has been explicitly classified. The summary's phrase ``rank-zero/linear part'' is undefined and is not adopted as a mathematical assertion.
\subsection{Short English statement}
Determine orbit closures in abelian- and quadratic-differential strata, their analytic or algebraic structure, and their geometric classification. The registry reports major structural results as solved, while its summary still treats a full classification as open; area-one finite-measure facts are preserved separately.
\subsection{Sources}
\begin{itemize}
\item Local UnsolvedMath registry, permanent ID 11000216, source code catalogue entry, Problem 9, ``Orbit closures for moduli spaces,'' PDF page 259. No individual proposer is supplied.
\item Farb, \emph{Problems on Mapping Class Groups and Related Topics}. Supplied URL: \url{https://www.math.uchicago.edu/~farb/papers/mcgbook.pdf}; recorded access 2026-07-29; extraction pdf-font-structure. The edited volume contains 329 explicitly typeset headings across chapters, with repeated labels disambiguated by page.
\item The source credits Kontsevich for affine period-coordinate structure and Masur--Veech for finite area-one measure. The summary credits Eskin--Mirzakhani--Mohammadi, Wright, Filip, and others for later work, using the phrase ``magic wand theorem.'' No full bibliographic entries or exact scope for the quadratic strata are supplied.
\item Supplied status solved and SOLVED-IN-LITERATURE; the summary says PARTIAL/SOLVED and distinguishes structural resolution from a still-active full classification. Generic background requires a current primary-literature audit and redistribution review, both NEEDS\_REVIEW. The registry assigns L4 whereas background records default L3.
\end{itemize}
\end{document}
